% Revised Results and Discussion working draft
% Short statements are intentional. Expand only where marked TODO.

\chapter{Results and Discussion}
\label{ch:results}

This chapter answers the three research questions. RQ1 compares top-height prediction under
spatial-block cross-validation, then looks at temporal transfer and a diagnostic check. RQ2
examines persistent plot-level departures from a shared Chapman--Richards (CR) curve. RQ3 examines
plot-specific asymptotic deviations, comparing global models to a spatially varying one (GNNWR).
Set~3 is used for the main RQ1 comparison as a common feature set across all model families, with
shared terrain and wind inputs. Full feature-set results are in Appendix~\ref{app:full-results}.

% Checked: Set3 first appears (as "nested_set3_gated_terrain_wind_vif") in the XGBoost
% hyperparameter search script, dated around the same time RQ1 results were being computed --
% no evidence it was fixed before test performance was inspected. Wording above uses the
% conservative "representative feature set" framing rather than a causal "because" justification.


%=======================================================================================
\section{RQ1: Prediction of top height}
\label{sec:results-rq1}

\subsection{Spatial-block performance}

Tuned XGBoost predicted top height best in both cohorts. It had the lowest RMSE and MAE, and the
highest $R^2$, under five-fold spatial-block cross-validation (Table~\ref{tab:results-rq1}). The
DNN was the strongest neural model. This directly answers the first part of RQ1: under realistic,
compartment-held-out testing, XGBoost beats every neural model tested here.

% Verified: Table~\ref{tab:results-rq1} matches the corrected/verified results table (Set~3,
% spatial\_block\_kfold, seed 42, tuned XGBoost from the 27-config search) exactly -- no
% discrepancy found against an earlier table.

\begin{table}[!htbp]
\centering
\scriptsize
\setlength{\tabcolsep}{3pt}
\renewcommand{\arraystretch}{0.95}
\begin{tabular}{lccc|ccc}
\toprule
& \multicolumn{3}{c|}{Four-survey}
& \multicolumn{3}{c}{Six-survey} \\
Model
& RMSE$\downarrow$ & MAE$\downarrow$ & $R^2\uparrow$
& RMSE$\downarrow$ & MAE$\downarrow$ & $R^2\uparrow$ \\
\midrule
Linear
& $5.168{\pm}0.361$ & $4.056{\pm}0.343$ & $0.580{\pm}0.026$
& $4.877{\pm}0.643$ & $3.619{\pm}0.426$ & $0.509{\pm}0.067$ \\
RF
& $4.794{\pm}0.187$ & $3.611{\pm}0.142$ & $0.638{\pm}0.011$
& $4.062{\pm}0.249$ & $3.010{\pm}0.185$ & $0.656{\pm}0.046$ \\
XGBoost (tuned)
& $\mathbf{4.548{\pm}0.137}$ & $\mathbf{3.431{\pm}0.123}$ & $\mathbf{0.674{\pm}0.009}$
& $\mathbf{3.656{\pm}0.232}$ & $\mathbf{2.669{\pm}0.183}$ & $\mathbf{0.722{\pm}0.031}$ \\
DNN
& $4.681{\pm}0.230$ & $3.533{\pm}0.201$ & $0.655{\pm}0.016$
& $3.903{\pm}0.316$ & $2.867{\pm}0.275$ & $0.684{\pm}0.031$ \\
PINN
& $5.206{\pm}0.399$ & $4.063{\pm}0.407$ & $0.573{\pm}0.036$
& $3.886{\pm}0.290$ & $2.819{\pm}0.196$ & $0.686{\pm}0.038$ \\
PINN-$k$
& $5.197{\pm}0.402$ & $4.046{\pm}0.421$ & $0.575{\pm}0.037$
& $3.895{\pm}0.282$ & $2.816{\pm}0.196$ & $0.684{\pm}0.038$ \\
\bottomrule
\end{tabular}
\caption{Five-fold spatial-block performance using Set~3. Values are the mean $\pm$ standard
deviation of the five fold-level metrics. Each compartment appears in one test fold. The pooled
held-out samples contain 232{,}292 plot--survey rows from 58{,}073 plots in the four-survey cohort
and 82{,}614 rows from 13{,}769 plots in the six-survey cohort.}
\label{tab:results-rq1}
\end{table}

The XGBoost lead is real, not an averaging artefact. RMSE was 2.8\% lower on 4survey and 6.3\%
lower on 6survey. A fold-by-fold check confirms it: XGBoost won 9 of the 10 individual folds
tested across both cohorts (all 5 on 6survey; 4 of 5 on 4survey, losing one fold by a small
margin, $-0.007$ $R^2$). This matches previous findings that boosted trees stay strong on
heterogeneous tabular data \citep{shwartzziv2021Tabular_C4_NN}.

Why does XGBoost win? Five candidate reasons were checked directly, not just assumed:

\begin{itemize}
\item \textbf{Feature interactions -- real, but not the answer.} 27\% of XGBoost's prediction
weight comes from interactions (SHAP check), mostly Age crossed with terrain and stand structure.
But the DNN already captures 80\% of the total possible nonlinear gain over a plain Linear model
(Table~\ref{tab:results-rq1}: $R^2$ 0.580 $\to$ 0.655 of a possible 0.580 $\to$ 0.674). Random
Forest -- also a nonlinear, interaction-capable tree model -- does worse than the DNN (0.638). So
the ability to model interactions is not what separates XGBoost from the DNN: the DNN already has
most of it. See `TEMP_rq1_xgb_interaction_check_2026-08-19.tex`.
\item \textbf{Too few independent compartments -- ruled out.} Subsampling 4survey down to
6survey's own compartment count (48, vs.\ 232 after filtering) did not widen XGBoost's lead. The
gap flipped sign and became noisy across random subsamples ($+0.024$ full vs.\ $-0.044$, $+0.029$,
$-0.085$ subsampled). See `TEMP_rq1_compartment_subsample_check_2026-08-19.tex`.
\item \textbf{DNN under-tuned or more hyperparameter-sensitive -- ruled out, on architecture.} An
architecture-size sweep found no consistent benefit for the DNN, and PINN/PINN$_k$ were even less
sensitive to architecture size, not more.
\item \textbf{Poor handling of unseen terrain in held-out compartments -- ruled out, mildly the
opposite.} Trees bound each leaf's prediction to what training rows looked like. A DNN's smooth
function can extrapolate further. Held-out compartments really can have terrain values never seen
in training, so this is a real problem worth checking, not just a guess. But XGBoost is not more
robust to it: on the 4.1\% of 4survey test rows with an out-of-range feature, XGBoost's error
rises MORE relative to its own normal error than the DNN's does (ratio $1.312$ vs.\ $1.243$), and
its error correlates slightly more, not less, with how far out of range a row is ($\rho=0.074$
vs.\ $0.059$). See `TEMP_rq1_extrapolation_check_2026-08-19.tex`.
\item \textbf{DNN under-tuned on training hyperparameters -- ruled out.} A 48-combination
dropout$\times$learning-rate grid search was run for all three neural models, both cohorts
(project-internal E7/E9 sweep), with the winning combination per model/cohort re-confirmed under
real pooled 5-fold cross-validation, not just the coarse single-split screen. On 4survey, every
winner's gain over the model's own default setting was within noise ($\pm0.008$ $R^2$); PINN's
4survey winner was identical to its default (dropout=0, lr=0.0001) -- the sweep itself found
nothing better. XGBoost's own 27-configuration grid search similarly moved its numbers by only
0.002--0.004 $R^2$ (Table~\ref{tab:results-rq1}'s tuned XGBoost row). Both models were searched;
neither was hiding a large hyperparameter-driven gain on 4survey.
\end{itemize}

Bottom line: five candidate reasons were checked. Four are ruled out: compartment count,
architecture size, extrapolation to unseen terrain, and training-hyperparameter tuning. One,
feature interactions, is real but does not decide it. The gap is real and consistent, but no
single tested mechanism explains it. The most likely remaining account is that XGBoost, as a
boosted-tree model, simply suits this tabular prediction task better than this DNN architecture
does -- consistent with the wider tabular-data literature already cited, not a fixable
methodological gap found in this chapter.

DNN also reliably beats PINN, but only on 4survey. Their 95\% intervals (1{,}000 cluster-bootstrap
resamples of whole compartments) do not overlap there (DNN $[0.624,0.690]$; PINN $[0.545,0.607]$).
On 6survey the intervals overlap heavily (DNN $[0.658,0.740]$; PINN $[0.671,0.736]$): a
statistical tie, not a real difference. So CR guidance does not just fail to beat XGBoost -- it
does not even reliably beat the plain DNN control on either cohort.


\subsection{Cohort and temporal sensitivity}

6survey is a different slice of the forest, not just a smaller sample of the same one. It has
fewer compartments (48 vs.\ 296), sits at lower and more uniform elevation ($102.3\pm47.9$m,
range 18--351m, vs.\ $177.9\pm104.6$m, range 18--561m), and its own fitted growth curve has a
lower ceiling but faster early growth ($y_{\max}=44.46$m, $k=0.0233$ vs.\ $y_{\max}=51.96$m,
$k=0.0103$). This composition difference, not sample size alone, is the more likely explanation
for 6survey's different behaviour throughout this chapter.

Every model lost accuracy under the 2023 temporal test, compared to spatial-block testing
(Table~\ref{tab:results-rq1-temporal}). This is expected: predicting further ahead in time is
harder than predicting a held-out compartment in the same period.

\begin{table}[!htbp]
\centering
\small
\renewcommand{\arraystretch}{1.2}
\setlength{\tabcolsep}{4pt}
\begin{tabular}{llrrr}
\toprule
Cohort & Model & Spatial-block $R^2$ & Temporal $R^2$ & Absolute change \\
\midrule
4survey & DNN & 0.634 & 0.544 & $-0.090$ \\
4survey & PINN & 0.584 & 0.290 & $-0.294$ \\
4survey & PINN$_k$ & 0.588 & 0.289 & $-0.299$ \\
6survey & DNN & 0.722 & 0.317 & $-0.405$ \\
6survey & PINN & 0.731 & 0.161 & $-0.570$ \\
6survey & PINN$_k$ & 0.730 & 0.190 & $-0.540$ \\
\bottomrule
\end{tabular}
\caption{Spatial-block vs.\ temporal-holdout $R^2$, by model and cohort. Training uses survey data
up to 2012, spatial-block validation uses 2021 data, temporal holdout tests on 2023 data. Only
$R^2$ was computed for this check; RMSE and MAE are not available from this source
\fc{(re-run if RMSE/MAE are needed for the final draft)}.}
\label{tab:results-rq1-temporal}
\end{table}

Physics guidance made the temporal drop worse, not better. PINN and PINN$_k$ lost far more
accuracy than the plain DNN, on both cohorts (4survey: DNN $-0.090$ vs.\ PINN $-0.294$; 6survey:
DNN $-0.405$ vs.\ PINN $-0.570$). If the physics constraint helped the model extrapolate, the drop
should be smaller for PINN, not larger. It is the opposite. This is a second respect, after
spatial-block accuracy, where CR guidance does not deliver its intended benefit.

Why does 6survey drop more than 4survey under temporal testing? Not the calendar gap: training
stops in 2012 and testing is 2023 for both cohorts alike, so that gap is identical for both.
Composition is the more likely reason: 6survey's narrower, lower-elevation sample gives every
model less to generalise from.


\subsection{Effect of the evaluation split}

Switching from a spatial-block split to a random (plot-level) split inflates the DNN's score
sharply: 0.831 vs.\ 0.634 on 4survey, 0.844 vs.\ 0.722 on 6survey. Nearby plots share terrain and
stand conditions, so a random split leaks information from train into test. Random-split scores
therefore overstate real-world accuracy -- this is why spatial-block splitting, not random
splitting, is used as the primary evaluation throughout this chapter.

PINN does not show this inflation: its random-split and spatial-split scores are nearly identical
(4survey: $0.590$ vs.\ $0.584$; 6survey: $0.728$ vs.\ $0.731$). This could come from the physics
constraint, from PINN's weaker overall fit, or from another architecture difference -- these were
not tested apart, so no single cause is confirmed.


\subsection{Physics-weight ablation}

Turning up the physics/trajectory loss weight made predictions worse, every time: both models,
both cohorts, all three weights tested (0, 1, 2). On 4survey this is a real effect, not noise:
removing the physics loss entirely ($\lambda=0$) raised $R^2$ by $0.05$ for both PINN and
PINN$_k$, outside their own bootstrap intervals. On 6survey the effect was smaller and within
sampling error ($+0.010$).

This is not just "no benefit" -- turning the weight up makes the core prediction actively worse,
at every step. That is direct evidence that the physics and trajectory loss terms compete with the
data-fitting objective on the same shared network weights, rather than acting as a neutral
regulariser. Removing them narrows most, but not all, of the DNN--PINN gap on 4survey (DNN
$R^2=0.655$, unconstrained PINN$_k$ $0.629$, down from a $0.080$ gap at $\lambda=1$). On 6survey it
closes the gap completely -- the unconstrained model even beats the DNN slightly ($0.719$ vs.\
$0.684$).

The gap narrows at $\lambda=0$, but does not close, and that residual gap is expected, not
puzzling. PINN's terrain and wind features never reach its height prediction directly -- they pass
only through a narrow sub-network used solely inside the physics/trajectory loss. At $\lambda=0$
that loss has zero weight, so the sub-network gets no training signal at all: the unconstrained
model is effectively a no-environment network, not an equivalent of the DNN, which feeds terrain
straight into its main trunk. So PINN's weakness has two separate, both now evidenced, causes: an
architectural bottleneck (terrain only ever reaches predictions through that narrow sub-network),
and loss competition (the physics/trajectory terms actively pull against accuracy once weighted
above zero).

Physics guidance does change the learned CR parameters. But the evidence that this is a real
benefit is incomplete. Without guidance ($\lambda=0$), $y_{\max}$ and $k$ collapse to just two
values per fold, not genuine per-plot values. So their reported correlation ($r=-0.935$) looks
precise but is not a real measure of per-plot identifiability. With guidance ($\lambda=1$) the
correlation is smaller ($r=-0.449$), but nobody has yet checked whether these values are genuinely
per-plot, or still partly collapsed. For context: the classical, non-neural CR fit has its own
correlation, $r=-0.993$, unaffected by this issue -- it supports the idea that some constraint is
needed. Bottom line: guidance changes the reported correlation. It does not yet prove that it
fixes identifiability.

% TODO: Confirm whether the lambda=1 (guided) learned y_max/k values vary genuinely per plot or
% also partially collapse -- the same raw-value check already done for lambda=0 has not been run
% for lambda=1, so r=-0.449's reliability is not yet established.
% Checked: no invalid lambda=0 CI/bootstrap number is currently stated in this chapter (the
% -0.935 point estimate above is reported honestly, with its precision caveat) -- nothing to move
% out.

% FIGURE 1 SPECIFICATION
% Chart: two-panel paired fold plot.
% Panel A: four-survey; Panel B: six-survey.
% x-axis: model; y-axis: fold R2. Join values from the same fold with thin grey lines.
% Highlight XGBoost, DNN and PINN; mute other baselines.
% Why: shows whether the headline ranking holds in each spatial fold. Adds information not visible
% in the mean+-SD table. Avoid another bar chart of the same averages.
\begin{figure}[!htbp]
\centering
\fbox{\parbox{0.92\textwidth}{\centering\vspace{3.0cm}
TODO FIGURE: paired fold-level $R^2$ comparison for XGBoost, DNN and PINN; one panel per cohort
\vspace{3.0cm}}}
\caption{Fold-level spatial performance of the main competing models. Thin lines connect models
evaluated on the same held-out compartment fold. This shows the consistency of the model ranking,
rather than repeating the table means.}
\label{fig:rq1-fold-performance}
\end{figure}

% FIGURE 2 SPECIFICATION
% Chart: line plot of lambda against held-out R2.
% Separate lines for PINN and PINN-k; panels for cohorts.
% Optional lower panel: a valid parameter-stability measure for lambda > 0 only.
% Why: directly shows the accuracy--constraint trade-off and whether it is monotonic.
% Do not plot the invalid lambda=0 parameter correlation as a plot-level estimate.

\paragraph*{Answer to RQ1.}
Tuned XGBoost gave the strongest spatially held-out predictions in both cohorts. The DNN remained
competitive, but CR guidance did not improve predictive accuracy. Random plot splitting gave an
optimistic DNN estimate. The current $\lambda=0$ implementation does not support a controlled
claim about parameter identifiability.


%=======================================================================================
\subsection{Diagnostic: where environmental models improve on the shared curve}
\label{sec:results-rq1-cr-reduction}

Environmental models trade accuracy across plots, they do not improve every plot equally: they
help most where the shared curve already fits worst, and often hurt where it already fits best.

\begin{table}[!htbp]
\centering
\scriptsize
\renewcommand{\arraystretch}{1.1}
\begin{tabular}{llrrrr}
\toprule
Model & Cohort & CR MAE (m) & Reduction (m) & Reduction (\%) & Rows improved (\%) \\
\midrule
DNN, environment & Four-survey & 5.028 & $1.501\pm0.023$ & 29.9 & $64.8\pm4.1$ \\
DNN, environment & Six-survey & 3.047 & $0.174\pm0.036$ & 5.7 & $52.4\pm4.8$ \\
XGBoost, environment & Four-survey & 5.028 & $1.597\pm0.330$ & 31.8 & $67.2\pm3.8$ \\
XGBoost, environment & Six-survey & 3.047 & $0.361\pm0.183$ & 11.8 & $56.0\pm4.6$ \\
XGBoost, no environment & Four-survey & 5.030 & $1.222\pm0.232$ & 24.3 & $64.5\pm3.2$ \\
XGBoost, no environment & Six-survey & 3.047 & $0.240\pm0.144$ & 7.9 & $54.8\pm3.0$ \\
\bottomrule
\end{tabular}
\caption{Reduction in absolute error relative to the fold-matched CR baseline. Positive values
indicate lower error than CR. The DNN uncertainty summarises five seeds; XGBoost uses one fixed
configuration. These uncertainty terms are therefore not directly equivalent.}
\label{tab:results-cr-reduction}
\end{table}

On 4survey, XGBoost cuts Q4 (worst-fit) error by 5.08~m but worsens Q1 (best-fit) error
$2.6\times$; 6survey shows the same pattern (Q4: $-1.81$m; Q1: $3.2\times$ worse). A model that is
only exploiting "more room to improve" in Q4 would show this pattern too, so this alone does not
prove real environmental correction. A no-environment XGBoost control shows real signal is still
part of the story: adding environmental inputs deepens the reduction by a further 0.375~m
(4survey) and 0.121~m (6survey) on top of the no-environment control.

A permutation test settles the "is this real or just room to improve" question directly. Feeding
XGBoost scrambled, meaningless environmental values still degrades Q1 by almost the same amount as
real values do (4survey: $-1.304$m fake vs.\ $-1.207$m real; 6survey: $-0.565$m vs.\ $-0.723$m).
So Q1's harm is not a real cost of environmental information -- it is the generic cost of giving
any flexible model room to vary its output on plots the baseline already fits well. Q4 is
different: fake values recover only 76--82\% of the real gain (4survey: 3.973m fake vs.\ 4.833m
real, an 18\% shortfall; 6survey: 1.180m vs.\ 1.614m, 27\%), so real environmental signal adds a
genuine benefit there, on top of flexibility alone. (DNN and PINN were not re-tested this way.)

This answers a key part of RQ1's "why": environmental conditioning genuinely helps the worst-fit
plots, using real signal, not noise. It has a real, non-artefactual cost on already-well-fit
plots too. Neither effect is just "more room to improve."

No stable spatial pattern explains where conditioning helps. 6survey shows no spatial clustering
at all. 4survey's apparent pattern disappears and flips sign at a wider search window -- the
signature of noise, not a real geographic pattern. A spatially organised "help map" is not
supported by this evidence.

% Move exact windows, sign changes and the 47-compartment subsampling check to the appendix.

% FIGURE 3 SPECIFICATION
% Chart: dot-and-line interaction plot, not grouped bars.
% x-axis: CR-error quartile Q1--Q4; y-axis: mean absolute-error reduction in metres.
% Panels: four-survey and six-survey.
% Lines: environment model, no-environment control, shuffled-environment control.
% Add a horizontal zero line. Positive = improvement; negative = harm.
% Why: shows the central trade-off, model controls and cohort replication in one figure. The table
% only gives the aggregate and misses this mechanism.
\begin{figure}[!htbp]
\centering
\fbox{\parbox{0.92\textwidth}{\centering\vspace{3.0cm}
TODO FIGURE: CR-error quartile against mean error reduction; controls and cohorts shown separately
\vspace{3.0cm}}}
\caption{Error reduction relative to CR across quartiles of CR baseline error. The zero line
separates improvement from degradation. This figure shows whether gains are concentrated among
observations that depart most from the shared curve.}
\label{fig:cr-error-quartiles}
\end{figure}


%=======================================================================================
\section{RQ2: Persistent departure from the shared curve}
\label{sec:results-rq2}

RQ2 models each plot's mean residual across its four-survey history. Canopy cover and thinning
variables had the largest attribution values across the tested feature sets and models.

\begin{table}[!htbp]
\centering
\scriptsize
\renewcommand{\arraystretch}{1.1}
\begin{tabular}{lcccc}
\toprule
Set & Elastic Net $R^2$ [95\% CI] & XGBoost $R^2$ [95\% CI]
& EN Moran's $I$ & XGB Moran's $I$ \\
\midrule
Set1 & 0.220 [0.180, 0.252] & 0.222 [0.183, 0.252] & 0.145 & 0.146 \\
Set2 & 0.359 [0.310, 0.396] & 0.416 [0.355, 0.461] & 0.105 & 0.032 \\
Set3 & 0.325 [0.272, 0.365] & 0.375 [0.321, 0.414] & 0.110 & 0.077 \\
Set4 & 0.358 [0.305, 0.398] & 0.395 [0.338, 0.438] & 0.103 & 0.054 \\
\bottomrule
\end{tabular}
\caption{Held-out attribution of each plot's mean CR residual in the four-survey cohort. Moran's
$I$ is evaluated using each model's semivariogram-resolved distance band. Full ranges, permutation
tests and mixed-effects results are reported in Appendix~\ref{app:full-results}.}
\label{tab:results-rq2}
\end{table}

Adding environmental features improved attribution over the structural baseline: every feature
set clears Set~1 (baseline-only) with no overlap. Set~2 had the highest XGBoost score. Its
interval overlaps Set~4, so we cannot say one clearly beats the other. Set~3 (terrain-only, no
soil or climate) scored lowest of the three. So soil and climate variables add real signal beyond
terrain and wind alone -- terrain alone is not enough.

Predictors were standardised (mean 0, SD 1) before fitting Elastic Net and the mixed-effects
model. So their coefficients mean: metres of CR residual per one-SD change in that predictor.
XGBoost does not need standardising, since it is tree-based. Its SHAP values sit on the original
response scale, so they are not directly comparable to the Elastic Net coefficients above.

Canopy cover had the largest Elastic Net coefficient and the largest XGBoost SHAP value. Removing
it from Set~4 lowered held-out $R^2$ in both models \sr{(exact removal-$R^2$ values not added: no
ablation script matching Set~4's folds and final model version was found this session -- verify
before citing a number)}. Canopy cover also correlates moderately with the height response. This
may reflect stand structure, past growth, and the fact that canopy cover and height share the same
LiDAR source. So canopy cover is treated as a structural control, not an independent environmental
driver.

Slope was the only abiotic predictor with a stable direction and a real size, across the main
models. TOPEX was positive in the mixed-effects model, but weak and unstable in Elastic Net. We
checked why, rather than guessing. TOPEX's variance is mostly between compartments: 74\%, versus
63\% for slope. The mixed-effects model has a compartment random intercept that can absorb this.
Elastic Net does not. So unmeasured compartment-level factors have more room to leak into TOPEX's
coefficient and shift it between folds. This is only part of the story: slope is also mostly
between-compartment (63\%), just less so -- the 11-point gap is real but modest on its own.
TOPEX's true effect also looks weaker and noisier than slope's: slope's coefficient size matches
across methods, while TOPEX's Elastic Net estimate is small relative to its own spread. A weak
true signal flips sign more easily under the same confounding pressure than a strong one does.
Together, these two checked facts plausibly explain the instability. But how they combine has not
itself been tested directly, so this is a well-supported account, not a fully proven mechanism.

As environmental features were added, XGBoost left less spatial pattern in its residuals than
Elastic Net did, though both start from a similar level. Moran's $I$ drops from 0.146 to 0.054 for
XGBoost, a 63\% fall. For Elastic Net it only drops from 0.145 to 0.103, a 29\% fall. \ai{Plausible
explanation: Elastic Net is linear, so it can only remove the linear part of each predictor's
spatial signal. Any nonlinear or interaction effect (like a terrain/wind threshold) stays in its
residual. Terrain predictors are themselves spatially smooth, so that leftover signal still
clusters by location. XGBoost's tree structure can fit these effects directly, removing more of
the spatial signal. This is the same model-flexibility gap RQ1's XGBoost-vs-DNN result raises -- a
hypothesis that fits the pattern, not a tested one.} \lim{Adding interaction or polynomial terms to
Elastic Net (or fitting a GAM with the same predictors) would directly test this, rather than
leaving it untested.}

% Checked: the second-stage model uses statsmodels MixedLM (models/spatial_attribution/nlme.py),
% a LINEAR mixed model fit to CR residuals (fixed terrain/wind effects, random compartment
% intercept) -- not a genuinely nonlinear mixed-effects fit. "mixed-effects model" (used
% throughout this section) is the correct label; "NLME" should not be used for it.
% TODO: Define the two-stage mixed-effects output in the appendix (fixed/random effect structure,
% what "mixed-effects model" refers to here) -- exact appendix location not established this
% session, content above is ready to use.

% FIGURE 4 SPECIFICATION
% Composite with two connected panels.
% Panel A: coefficient/SHAP rank plot for only the 6--8 most important features. Show direction for
% Elastic Net and mixed effects; magnitude for SHAP. Mark predictors with stable sign across folds.
% Panel B: residual Moran's I by Set for EN and XGBoost, with connected points.
% Why: Panel A answers what explains departure. Panel B shows what remains spatially unexplained.
% A residual map may be a small inset only if it reveals a pattern not conveyed by Moran's I.

\paragraph*{Answer to RQ2.}
Stand structure and management explained more persistent CR departure than any single abiotic
predictor. Slope was the most stable abiotic association. Environmental features improved held-out
attribution, but residual spatial autocorrelation remained.


%=======================================================================================
\section{RQ3: Plot-specific asymptotic deviation}
\label{sec:results-rq3}

RQ3 examines the difference between each plot's fitted asymptote and its yield-class expectation.
It compares global attribution models with GNNWR, which allows relationships to vary by location.

\begin{table}[!htbp]
\centering
\scriptsize
\renewcommand{\arraystretch}{1.1}
\begin{tabular}{lcc|cc|cc}
\toprule
& \multicolumn{2}{c|}{Elastic Net $R^2$}
& \multicolumn{2}{c|}{XGBoost $R^2$}
& \multicolumn{2}{c}{GNNWR $R^2$} \\
Set & Four-survey & Six-survey & Four-survey & Six-survey & Four-survey & Six-survey \\
\midrule
Set2 & 0.286 [0.241, 0.325] & 0.057 & 0.266 [0.214, 0.308] & 0.031
& 0.320 [0.270, 0.369] & $0.014\pm0.053$ \\
Set3 & 0.274 [0.223, 0.316] & 0.031 & 0.281 [0.235, 0.320] & 0.086
& 0.319 [0.263, 0.365] & $0.054\pm0.055$ \\
Set4 & 0.240 [0.192, 0.286] & 0.056 & 0.250 [0.201, 0.295] & 0.015
& 0.294 [0.236, 0.347] & $0.002\pm0.068$ \\
\bottomrule
\end{tabular}
\caption{Held-out $R^2$ for plot-specific asymptotic deviation. Four-survey brackets are 95\%
cluster-bootstrap intervals. Six-survey GNNWR values are the mean $\pm$ standard deviation across
three split seeds. Other six-survey entries are single estimates. These uncertainty summaries are
not directly equivalent.}
\label{tab:results-rq3}
\end{table}

GNNWR scored highest on 4survey, in all three feature sets. But its confidence intervals overlap
the global models' -- so this suggests an advantage, it does not prove one. GNNWR also left less
spatial pattern in its residuals (lower Moran's $I$) in every 4survey feature set.

All models did poorly on 6survey. No stable spatial pattern showed up in the residuals, at any
distance tested -- so Moran's $I$ cannot explain the near-zero accuracy. Subsampling 4survey down
to 6survey's own compartment count did not reproduce this behaviour either. So cohort composition,
not just compartment count, is the more likely explanation.

% TODO: move the block below into Appendix~\ref{app:full-results} (exact appendix file/location
% not established in this session -- content is ready, placement is not).
% Appendix content -- exact semivariogram ranges, p-values, 10km windows, subsample values:
%   RQ2 (mean CR residual, 4survey): ranges 1665-1739m (EN), 1385-1675m (XGB), all p=0.001.
%   RQ3 (asymptote deviation): 4survey ranges 808-1252m, all p=0.001. 6survey: Set3 unresolved
%     (flat semivariogram) for EN/XGB/GNNWR; Set2/Set4 XGB+GNNWR only resolve after widening to
%     10000m, landing at I~-0.0002, p=0.05-0.31 (not significant); 6survey EN unresolved in all
%     3 sets.
%   47-compartment subsample check (4survey, matching 6survey's compartment count): 3 random
%     seeds, all resolve cleanly, I=0.18-0.36, ranges 1045-1659m -- rules out compartment count
%     alone as the explanation for 6survey's unresolvability.

Canopy cover ranked highly in the 4survey global and spatial models. Other feature rankings
differed between methods -- read these cautiously, since each method represents spatial structure
differently. We do not interpret 6survey feature importance at all: held-out $R^2$ there is near
zero (Table~\ref{tab:results-rq3}, $0.014$--$0.086$), so any ranking there would just be noise.


\subsection{Extreme asymptotic deviations}

Tukey fences flagged 266 4survey plots (0.47\%) and 709 6survey plots (5.32\%). Their raw height
trajectories usually looked plausible, but some fitted asymptotes did not. Here is why: fixing the
growth-rate and shape parameters can push the fitted asymptote too high, whenever a plot grows
more slowly than its assigned yield-class curve expects. Standard raw-height filters miss this,
because it is a problem with how the target is built, not with the raw data.

Removing flagged plots did not improve accuracy. 4survey $R^2$ barely moved ($-0.002$). 6survey
accuracy got worse. So model accuracy alone does not justify excluding these plots. Even so, their
fitted asymptotes should not be read as real biological ceilings.

None of the 266 flagged 4survey plots met the clearfelling or measurement-anomaly rules
\fc{(exact rule definitions still need writing up in the methodology/appendix, not yet
documented)}. So active disturbance is not the cause. This supports the target-construction
explanation above: the flagged asymptotes are an artefact of how the target is built, not real
biological ceilings or unrecorded felling. Flagged plots sit closer to compartment boundaries than
the rest of the population (median 9.4~m vs.\ 28.0~m), and were more often poorly fitted in 2008
(75.6\% vs.\ 47.8\% baseline). This points to boundary and yield-class assignment error, more than
any single cause. Other possible factors: clipped plot geometry, boundary changes, mixed stand
conditions. None of these were tested apart from each other.

% FIGURE 5 SPECIFICATION
% Three-panel explanatory figure.
% A: map of asymptotic deviation, with flagged plots outlined rather than dominating the colour
% scale. Use a robust symmetric colour scale based on central quantiles.
% B: three example height--age trajectories: stable fit, moderate deviation, unstable high
% asymptote. Plot observed heights, fixed-shape fitted curve and yield-class reference.
% C: compact distribution of deviation by distance-to-boundary group or flagged status.
% Why: explains what the target means, why extremes occur, and where they occur. This is more
% informative than another model-score chart.
\begin{figure}[!htbp]
\centering
\fbox{\parbox{0.92\textwidth}{\centering\vspace{3.0cm}
TODO FIGURE: deviation map, three trajectory archetypes, and boundary comparison
\vspace{3.0cm}}}
\caption{Spatial distribution and construction of plot-specific asymptotic deviations. Example
trajectories distinguish plausible departures from high asymptotes produced by the fixed-shape
fit. Flagged plots remain in the main analysis.}
\label{fig:rq3-deviation-examples}
\end{figure}

% OPTIONAL FIGURE 6 SPECIFICATION
% Chart: slope chart or dot plot of residual Moran's I for EN, XGBoost and GNNWR across Sets 2--4.
% Four-survey only. Pair models within each set. Include semivariogram range in caption/table, not
% encoded in the same chart.
% Why: supports the spatial-modelling claim without repeating R2 values.

\paragraph*{Answer to RQ3.}
GNNWR scored highest on 4survey and left the least spatial pattern in its residuals. But
overlapping intervals mean this is not a proven advantage. No model generalised well on 6survey.
Extreme target values usually trace back to instability in how the fixed-shape asymptote is built
-- they are not real ecological ceilings.


%=======================================================================================
\section{Overall discussion}
\label{sec:results-overall-discussion}

Three findings tie the research questions together. First, tuned XGBoost gave the strongest
top-height predictions. Second, CR guidance lowered accuracy at every loss weight tested. Third,
stand structure and management explained more curve departure than any single abiotic predictor
-- though spatial structure was still left over.

PINN's weak result does not mean CR theory is wrong for forest growth. It means the tested
derivative and trajectory penalties did not improve prediction on this data. One reason: a single
pooled CR derivative may not fit local trajectories that differ due to management, disturbance, or
site conditions. Another: the DNN may already capture most of the predictable age--environment
relationship without needing the extra penalty.

The attribution results show association, not causation. The data come from one forest only.
Several environmental variables are spatially coarse or model-derived, not measured directly.
Management records are incomplete. Canopy cover shares its LiDAR source with the height response
itself. These limits restrict how far the ecological interpretation can be trusted, and whether it
transfers beyond Aberfoyle.

% TODO: Expand by one paragraph only:
% - practical meaning for forest monitoring;
% - why spatial validation changes the conclusion;
% - what a future external-site test should evaluate.
% Keep detailed future experiments for the Conclusion/Future Work chapter.

